\newpage
\section{Exhaustive Practice Problem Bank (Chapter 2)}
\label{sec:ch02_problem_bank}

\begin{tcolorbox}[enhanced,colback=subtlegreen,colframe=cengagegreen,arc=2mm,boxrule=1pt,title=\textbf{\large Organization of Practice Problems}]
All questions are organized by \textbf{Sub-topic} with full origin book citations. In the Complete Solutions Edition, full mathematical and chemical derivations are provided directly beneath each problem.
\end{tcolorbox}

% ==============================================================================
% SUB-TOPIC 2.1: CELL FUNDAMENTALS, POLARITY & IUPAC REPRESENTATION
% ==============================================================================
\subsection{Sub-Topic 2.1: Cell Fundamentals, Polarity \& IUPAC Representation}

% Problem 2.1
\begin{problembox}
\textbf{Problem 2.1} \hfill \textbf{[Source: Narendra Avasthi Level-1 Q1]}

A cell redox reaction would be thermodynamically spontaneous if the standard cell potential ($E^\circ_{\text{cell}}$) and standard Gibbs free energy change ($\Delta_r G^\circ$) are respectively:
\begin{multicols}{2}
\begin{enumerate}[label=(\alph*)]
    \item Positive and negative
    \item Negative and negative
    \item Zero and zero
    \item Positive and zero
\end{enumerate}
\end{multicols}
\end{problembox}
\begin{solution}
\textbf{Correct Option: (a)}

\textbf{Step-by-Step Derivation:}
The fundamental relation connecting electrical work and free energy change is:
\begin{equation*}
    \Delta_r G^\circ = -n F E^\circ_{\text{cell}}
\end{equation*}
For any process to be spontaneous under constant temperature and pressure, the second law of thermodynamics mandates $\Delta_r G < 0$. Therefore, since $n$ and $F$ are strictly positive constants:
\begin{equation*}
    \Delta_r G^\circ < 0 \iff E^\circ_{\text{cell}} > 0 \quad (\text{Positive})
\end{equation*}
Hence, $E^\circ_{\text{cell}}$ must be \textbf{positive} and $\Delta_r G^\circ$ must be \textbf{negative}.
\end{solution}

% Problem 2.2
\begin{problembox}
\textbf{Problem 2.2} \hfill \textbf{[Source: Narendra Avasthi Level-1 Q2]}

Which of the following statements regarding electrode polarities is correct?
\begin{enumerate}[label=(\alph*)]
    \item Cathode is the negative terminal in both galvanic and electrolytic cells.
    \item Anode is the positive terminal in both galvanic and electrolytic cells.
    \item Cathode is the negative terminal in an electrolytic cell, and positive in a galvanic cell.
    \item Cathode is the positive terminal in an electrolytic cell, and negative in a galvanic cell.
\end{enumerate}
\end{problembox}
\begin{solution}
\textbf{Correct Option: (c)}

\textbf{Step-by-Step Derivation:}
By fundamental definition:
\begin{itemize}
    \item \textbf{Anode:} Electrode where oxidation occurs.
    \item \textbf{Cathode:} Electrode where reduction occurs.
\end{itemize}
\begin{enumerate}[leftmargin=*]
    \item In a \textbf{Galvanic Cell}: Oxidation at the anode releases electrons, making it the negative terminal ($-$). The cathode consumes electrons, acquiring positive potential ($+$).
    \item In an \textbf{Electrolytic Cell}: An external battery pumps electrons into the cathode (making it negative, $-$) to force reduction, and pulls electrons from the anode (making it positive, $+$) to force oxidation.
\end{enumerate}
Thus, the cathode is \textbf{negative in an electrolytic cell and positive in a galvanic cell}.
\end{solution}

% Problem 2.3
\begin{problembox}
\textbf{Problem 2.3} \hfill \textbf{[Sources: Neeraj Kumar Ex-I Q1 / Pearson Ex-I Q1]}

A metal rod is immersed in an aqueous solution of its own ions. Its electrode potential is independent of:
\begin{multicols}{2}
\begin{enumerate}[label=(\alph*)]
    \item Temperature of the solution
    \item Concentration of the metal ions
    \item Area of the metal surface exposed
    \item Nature of the metal
\end{enumerate}
\end{multicols}
\end{problembox}
\begin{solution}
\textbf{Correct Option: (c)}

\textbf{Step-by-Step Derivation:}
By the Nernst equation for a metal-metal ion electrode:
\begin{equation*}
    E = E^\circ - \frac{RT}{zF} \ln \frac{1}{[\ce{M^{z+}}]}
\end{equation*}
The electrode potential is an \textbf{intensive property} determined solely by:
\begin{itemize}
    \item The intrinsic chemical nature of the metal ($E^\circ$).
    \item The concentration (activity) of the metal ions in solution ($[\ce{M^{z+}}]$).
    \item Temperature ($T$).
\end{itemize}
It is completely independent of the geometric dimensions or surface area of the electrode exposed to the electrolyte.
\end{solution}

% ==============================================================================
% SUB-TOPIC 2.2: SALT BRIDGE MECHANICS & REFERENCE ELECTRODES
% ==============================================================================
\subsection{Sub-Topic 2.2: Salt Bridge Mechanics \& Reference Electrodes}

% Problem 2.4
\begin{problembox}
\textbf{Problem 2.4} \hfill \textbf{[Sources: Pearson Ex-II Q14 / Neeraj Kumar Ex-II Q12]}

In a galvanic cell, the salt bridge:
\begin{enumerate}[label=(\alph*)]
    \item Does not participate chemically in the cell reaction.
    \item Stops the mutual bulk diffusion of ions from one electrode compartment to another.
    \item Is necessary for maintaining electrical neutrality and eliminating liquid junction potential.
    \item Ensures complete mechanical mixing of the two electrolytic solutions.
\end{enumerate}
\end{problembox}
\begin{solution}
\textbf{Correct Options: (a), (b), (c)}

\textbf{Step-by-Step Derivation:}
A salt bridge contains an inert electrolyte ($\ce{KCl}, \ce{KNO3}$) in agar-agar gel.
\begin{itemize}
    \item (a) is correct: The ions of the salt bridge do not take part in the redox reaction.
    \item (b) is correct: The gel matrix prevents the physical intermixing and bulk convection of the two half-cell solutions.
    \item (c) is correct: Ion migration through the bridge maintains electrical neutrality in both beakers and minimizes liquid junction potential.
    \item (d) is incorrect: It prevents mixing rather than ensuring it.
\end{itemize}
\end{solution}

% Problem 2.5
\begin{problembox}
\textbf{Problem 2.5} \hfill \textbf{[Source: Cengage Solved Ex 3.2]}

A $\ce{KCl}$ salt bridge cannot be used in which of the following electrochemical cells?
\begin{multicols}{2}
\begin{enumerate}[label=(\alph*)]
    \item $\ce{Zn | ZnSO4} \parallel \ce{CuSO4 | Cu}$
    \item $\ce{Fe | FeSO4} \parallel \ce{CuSO4 | Cu}$
    \item $\ce{Ag | AgNO3} \parallel \ce{CuSO4 | Cu}$
    \item $\ce{Ni | NiSO4} \parallel \ce{CuSO4 | Cu}$
\end{enumerate}
\end{multicols}
\end{problembox}
\begin{solution}
\textbf{Correct Option: (c)}

\textbf{Step-by-Step Derivation:}
When $\ce{KCl}$ is used in a cell containing $\ce{Ag+}$ ions ($\ce{AgNO3}$ half-cell), $\ce{Cl-}$ ions diffusing from the salt bridge react immediately with $\ce{Ag+}$ to form an insoluble precipitate of silver chloride:
\begin{equation*}
    \ce{Ag+(aq) + Cl^-(aq) -> AgCl(s) v}
\end{equation*}
The precipitation of $\ce{AgCl(s)}$ blocks the porous plugs of the salt bridge, disrupts ionic conduction, and alters the $[\ce{Ag+}]$ activity. Therefore, a $\ce{KNO3}$ or $\ce{NH4NO3}$ salt bridge must be used instead.
\end{solution}

% Problem 2.6
\begin{problembox}
\textbf{Problem 2.6} \hfill \textbf{[Source: Neeraj Kumar Ex-I Q8]}

The potential of a hydrogen electrode placed in a solution of $\text{pH} = 10$ at $298\text{ K}$ and $P_{\ce{H2}} = 1\text{ atm}$ is:
\begin{multicols}{2}
\begin{enumerate}[label=(\alph*)]
    \item $+0.591\text{ V}$
    \item $-0.591\text{ V}$
    \item $0.00\text{ V}$
    \item $-0.0591\text{ V}$
\end{enumerate}
\end{multicols}
\end{problembox}
\begin{solution}
\textbf{Correct Option: (b)}

\textbf{Step-by-Step Derivation:}
For the reduction half-reaction: $\ce{2H+(aq) + 2e^- <=> H2(g)}$.
\begin{equation*}
    E = E^\circ - \frac{0.0591}{2} \log \frac{P_{\ce{H2}}}{[\ce{H+}]^2}
\end{equation*}
Since $E^\circ = 0.00\text{ V}$ and $P_{\ce{H2}} = 1\text{ atm}$:
\begin{equation*}
    E = -0.0591 \times (-\log [\ce{H+}]) = -0.0591 \times \text{pH}
\end{equation*}
Given $\text{pH} = 10$:
\begin{equation*}
    E = -0.0591 \times 10 = -0.591\text{ V}
\end{equation*}
\end{solution}

% ==============================================================================
% SUB-TOPIC 2.3: ELECTROCHEMICAL SERIES & DISPLACEMENT PHENOMENA
% ==============================================================================
\subsection{Sub-Topic 2.3: Electrochemical Series \& Displacement Phenomena}

% Problem 2.7
\begin{problembox}
\textbf{Problem 2.7} \hfill \textbf{[Source: Pearson Ex-II Section A Q2]}

A student made the following experimental observations in the laboratory:
\begin{enumerate}[label=(\Roman*)]
    \item Clean copper metal did not react when immersed in $1.0\text{ M } \ce{Pb(NO3)2}$ solution.
    \item Clean lead metal dissolved when placed in $1.0\text{ M } \ce{AgNO3}$ solution and lustrous crystals of $\ce{Ag}$ metal deposited on its surface.
    \item Clean silver metal did not react when immersed in $1.0\text{ M } \ce{Cu(NO3)2}$ solution.
\end{enumerate}
The correct order of decreasing reducing character of the three metals is:
\begin{multicols}{2}
\begin{enumerate}[label=(\alph*)]
    \item $\ce{Cu} > \ce{Pb} > \ce{Ag}$
    \item $\ce{Ag} > \ce{Cu} > \ce{Pb}$
    \item $\ce{Pb} > \ce{Cu} > \ce{Ag}$
    \item $\ce{Pb} > \ce{Ag} > \ce{Cu}$
\end{enumerate}
\end{multicols}
\end{problembox}
\begin{solution}
\textbf{Correct Option: (c)}

\textbf{Step-by-Step Derivation:}
\begin{itemize}
    \item Observation (I): $\ce{Cu}$ cannot displace $\ce{Pb^2+} \implies \ce{Pb}$ is more electropositive than $\ce{Cu}$, meaning $\ce{Pb}$ is a stronger reducing agent than $\ce{Cu}$ ($\ce{Pb} > \ce{Cu}$).
    \item Observation (II): $\ce{Pb}$ displaces $\ce{Ag+} \implies \ce{Pb}$ is a stronger reducing agent than $\ce{Ag}$ ($\ce{Pb} > \ce{Ag}$).
    \item Observation (III): $\ce{Ag}$ cannot displace $\ce{Cu^2+} \implies \ce{Cu}$ is a stronger reducing agent than $\ce{Ag}$ ($\ce{Cu} > \ce{Ag}$).
\end{itemize}
Combining the three observations gives the unambiguous sequence:
\begin{equation*}
    \ce{Pb} > \ce{Cu} > \ce{Ag}
\end{equation*}
\end{solution}

% Problem 2.8
\begin{problembox}
\textbf{Problem 2.8} \hfill \textbf{[Source: Pearson Ex-II Q12]}

For the reduction of $\ce{NO3-}$ ion in an acidic aqueous solution:
\begin{equation*}
    \ce{NO3- + 4H+ + 3e^- -> NO(g) + 2H2O} \quad (E^\circ = +0.96\text{ V})
\end{equation*}
Given the standard reduction potentials of some metal couples:
\begin{align*}
    \ce{V^2+ + 2e^- <=> V(s)} \quad & E^\circ = -1.19\text{ V} \\
    \ce{Fe^3+ + 3e^- <=> Fe(s)} \quad & E^\circ = -0.04\text{ V} \\
    \ce{Hg2^2+ + 2e^- <=> 2Hg(l)} \quad & E^\circ = +0.86\text{ V} \\
    \ce{Au^3+ + 3e^- <=> Au(s)} \quad & E^\circ = +1.40\text{ V}
\end{align*}
Which of the following metals can be oxidized by $\ce{NO3-}$ in acidic aqueous solution under standard conditions?
\begin{multicols}{2}
\begin{enumerate}[label=(\alph*)]
    \item $\ce{V}$ and $\ce{Hg}$
    \item $\ce{Hg}$ and $\ce{Fe}$
    \item $\ce{Fe}$ and $\ce{V}$
    \item All except $\ce{Au}$
\end{enumerate}
\end{multicols}
\end{problembox}
\begin{solution}
\textbf{Correct Option: (d)}

\textbf{Step-by-Step Derivation:}
For $\ce{NO3-}$ to oxidize a metal $\ce{M}$ to its cation, the overall cell potential must be positive ($E^\circ_{\text{cell}} > 0$):
\begin{equation*}
    E^\circ_{\text{cell}} = E^\circ(\ce{NO3-/NO}) - E^\circ(\ce{M^{z+}/M}) = +0.96\text{ V} - E^\circ(\ce{M^{z+}/M}) > 0
\end{equation*}
\begin{equation*}
    E^\circ(\ce{M^{z+}/M}) < +0.96\text{ V}
\end{equation*}
Comparing the values:
\begin{itemize}
    \item For $\ce{V}$: $-1.19\text{ V} < +0.96\text{ V} \implies E^\circ_{\text{cell}} = +0.96 - (-1.19) = +2.15\text{ V} > 0$ (Oxidized).
    \item For $\ce{Fe}$: $-0.04\text{ V} < +0.96\text{ V} \implies E^\circ_{\text{cell}} = +0.96 - (-0.04) = +1.00\text{ V} > 0$ (Oxidized).
    \item For $\ce{Hg}$: $+0.86\text{ V} < +0.96\text{ V} \implies E^\circ_{\text{cell}} = +0.96 - (+0.86) = +0.10\text{ V} > 0$ (Oxidized).
    \item For $\ce{Au}$: $+1.40\text{ V} > +0.96\text{ V} \implies E^\circ_{\text{cell}} = +0.96 - (+1.40) = -0.44\text{ V} < 0$ (Cannot be oxidized).
\end{itemize}
Therefore, $\ce{V}, \ce{Fe},$ and $\ce{Hg}$ are oxidized by nitric acid under standard conditions, while gold ($\ce{Au}$) is immune.
\end{solution}

% Problem 2.9
\begin{problembox}
\textbf{Problem 2.9 (Arbitrary Shift of Reference Scale)} \hfill \textbf{[Source: Pearson Ex-II Comprehension I]}

Suppose that the Standard Hydrogen Electrode (S.H.E.) was arbitrarily assigned a standard reduction potential of $+1.00\text{ V}$ instead of $0.00\text{ V}$.
Given that on the standard scale:
$E^\circ(\ce{Zn^2+/Zn}) = -0.76\text{ V}$ and $E^\circ(\ce{Cu^2+/Cu}) = +0.34\text{ V}$.

\begin{enumerate}[label=\textbf{Part \arabic*:}]
    \item What would be the new observed standard electrode potential of the $\ce{Cu^2+/Cu}$ couple?
    \begin{multicols}{2}
    \begin{enumerate}[label=(\alph*)]
        \item $+0.34\text{ V}$
        \item $+1.34\text{ V}$
        \item $-0.66\text{ V}$
        \item $+0.66\text{ V}$
    \end{enumerate}
    \end{multicols}
    
    \item What would be the new standard EMF of the complete $\ce{Zn - Cu}$ Daniell cell?
    \begin{multicols}{2}
    \begin{enumerate}[label=(\alph*)]
        \item $+1.10\text{ V}$
        \item $+2.10\text{ V}$
        \item $+0.10\text{ V}$
        \item Indeterminate
    \end{enumerate}
    \end{multicols}
\end{enumerate}
\end{problembox}
\begin{solution}
\textbf{Part 1: Correct Option: (b)} \\
\textbf{Part 2: Correct Option: (a)}

\textbf{Step-by-Step Derivation:}
\textbf{Part 1: New Electrode Potential of $\ce{Cu}$:}
When the reference datum is shifted upward by $+1.00\text{ V}$, the absolute difference between any electrode and the SHE remains invariant:
\begin{equation*}
    E^\circ_{\text{new}}(\ce{M}) = E^\circ_{\text{old}}(\ce{M}) + 1.00\text{ V}
\end{equation*}
For copper:
\begin{equation*}
    E^\circ_{\text{new}}(\ce{Cu^2+/Cu}) = +0.34\text{ V} + 1.00\text{ V} = +1.34\text{ V}
\end{equation*}

\textbf{Part 2: Cell EMF of Daniell Cell:}
For zinc:
\begin{equation*}
    E^\circ_{\text{new}}(\ce{Zn^2+/Zn}) = -0.76\text{ V} + 1.00\text{ V} = +0.24\text{ V}
\end{equation*}
The cell EMF is the potential difference between the two electrodes:
\begin{equation*}
    E^\circ_{\text{cell}} = E^\circ_{\text{new}}(\text{cathode}) - E^\circ_{\text{new}}(\text{anode}) = (+1.34\text{ V}) - (+0.24\text{ V}) = +1.10\text{ V}
\end{equation*}
The cell EMF is a physically measurable, absolute voltage difference between the two terminal wires. It is \textbf{strictly invariant} under any arbitrary choice of reference zero.
\end{solution}

% ==============================================================================
% CHAPTER 2 ANSWER KEY TABLE (FOR MAIN.TEX)
% ==============================================================================
\ifshowsolutions
\else
\newpage
\subsection*{Answer Key: Chapter 2 Practice Problem Bank}
\addcontentsline{toc}{subsection}{Answer Key: Chapter 2}

\begin{center}
\renewcommand{\arraystretch}{1.3}
\begin{tabular}{|c|c||c|c||c|c|}
\hline
\textbf{Problem} & \textbf{Answer} & \textbf{Problem} & \textbf{Answer} & \textbf{Problem} & \textbf{Answer} \\
\hline
2.1 & (a) & 2.4 & (a), (b), (c) & 2.7 & (c) \\
2.2 & (c) & 2.5 & (c) & 2.8 & (d) \\
2.3 & (c) & 2.6 & (b) & 2.9 & Part 1: (b), Part 2: (a) \\
\hline
\end{tabular}
\end{center}
\fi
